% IEEE conference paper, on the IEEEtran class that comes with TeX Live.
% Guidelines: https://www.ieee.org/conferences/publishing/templates.html
\documentclass[conference]{IEEEtran}
\IEEEoverridecommandlockouts
\usepackage{cite}
\usepackage{amsmath,amssymb,amsfonts}
\usepackage{graphicx}
\usepackage{textcomp}
\usepackage{xcolor}

\begin{document}

\title{Paper Title}

\author{\IEEEauthorblockN{First Author}
\IEEEauthorblockA{\textit{Department} \\
\textit{Organization}\\
City, Country \\
email@example.com}
\and
\IEEEauthorblockN{Second Author}
\IEEEauthorblockA{\textit{Department} \\
\textit{Organization}\\
City, Country \\
email@example.com}
}

\maketitle

\begin{abstract}
A short summary of the problem, the approach and the results, in one paragraph.
\end{abstract}

\begin{IEEEkeywords}
first keyword, second keyword, third keyword
\end{IEEEkeywords}

\section{Introduction}
Say what the problem is and why it matters, then what this paper does~\cite{lamport1994}.

\section{Method}
\begin{equation}
  y = f(x) + \epsilon
  \label{eq:model}
\end{equation}
Equation~\eqref{eq:model} is referred to by its label.

\begin{figure}[htbp]
  \centering
  \includegraphics[width=\columnwidth]{example-image}
  \caption{A figure that spans one column.}
  \label{fig:example}
\end{figure}

\section{Experiments}
Results go here; see Fig.~\ref{fig:example}.

\section{Conclusion}
What was shown, and what comes next.

\section*{Acknowledgment}
Funding and people to thank.

\bibliographystyle{IEEEtran}
\bibliography{references}

\end{document}
